\section{Experiments}

Here, we present experiments quantizing the Llama family of models with \qts \cite{llama1,llama2,llama3}.
These models offer strong performance across a wide range of sizes, allowing us to compare how different quantization methods perform and scale.
We primarily compare \qts against \qss and AQLM.
For Llama 1, we include GPTVQ-2D instead of AQLM since AQLM does not publish Llama 1 numbers \cite{vanbaalen-gptvq}. 
GPTVQ-2D performs 2D VQ inside GPTQ and offers strong performance.
These methods outperform scalar quantization methods including GPTQ, AWQ, and OmniQuant; comparisons to those methods can be found in \qss and AQLM \cite{awq,optq,omniquant,tseng2024quip,aqlm}.
%Section \ref{sec:computedres} shows \qts with the 1MAD and 3INST computed codes, and section \ref{sec:hybres} shows how \qts performs with the hybrid lookup-computed code.
We mainly focus on the hybrid code (Section \ref{sec:hybres}) since it is tailored for modern GPUs, and present a full suite of results for it.
For the computed codes (Section \ref{sec:computedres}), we present results for Llama 2. 

Since the proxy error is not an additive distortion metric, we cannot minimize it by quantizing $W$ as one sequence.
Instead, for all experiments, we use \qts as a quantizer in \qs's BlockLDLQ, which allows us to simultaneously achieve high dimensionality and low proxy error \cite{tseng2024quip}.
Specifically, we quantize a block of $T_x \times T_y$ weights as a sequence, where $T_x$ and $T_y$ span the output and input dimensions of $W$, respectively.
Since BlockLDLQ only specifies feedback along the input dimension, this is equivalent to BlockLDLQ with $g=T_y$ but a vector dimension of $T_xT_y \gg T_y$. 
This has the benefit of limiting the effect of $g$ in BlockLDLQ's error bound $gm\mu^2\sigma^2\mbox{tr}(H^{1/2})^2/n$ while achieving a high dimension for TCQ.
Algorithm \ref{alg:qtblockldlq} in the Appendix describes this in more detail.
\subsection{Lookup-Free Computed Codes}
\label{sec:computedres}
% Please add the following required packages to your document preamble:
% \usepackage{multirow}
\begin{table}[t]
\caption{Wikitext2 and C4 perplexity ($\downarrow$), ctx. 4096, \qts with pure-computed codes. Even \textit{without fine-tuning}, pure-computed \qts outperforms \qss and AQLM, both of which use fine-tuning, at almost all models sizes.}
\centering
\small\sc
\tabcolsep=0.075cm
%\renewcommand{\arraystretch}{0.9}
\begin{tabular}{@{}cccccccccccccccccc@{}}
\multicolumn{1}{l}{}                           & \multicolumn{1}{l}{}                            & \multicolumn{1}{l}{}                              & \multicolumn{3}{c}{4 Bit No FT}                                                   & \multicolumn{2}{c}{$\approx$4 Bit}                        & \multicolumn{3}{c}{3 Bit No FT}                                                   & \multicolumn{2}{c}{$\approx$3 Bit}                        & \multicolumn{3}{c}{2 Bit No FT}                                                   & \multicolumn{2}{c}{$\approx$2 Bit} \\ \midrule
\multicolumn{1}{l}{}                           & \multicolumn{1}{l|}{}                           & \multicolumn{1}{c|}{\tiny FP16}                         & \tiny 1MAD          & \tiny 3INST         & \multicolumn{1}{c|}{\tiny QuIP\#}                       & \tiny QuIP\# & \multicolumn{1}{c|}{\tiny AQLM}                         & \tiny 1MAD          & \tiny 3INST         & \multicolumn{1}{c|}{\tiny QuIP\#}                       & \tiny QuIP\# & \multicolumn{1}{c|}{\tiny AQLM}                         & \tiny 1MAD          & \tiny 3INST         & \multicolumn{1}{c|}{\tiny QuIP\#}                       & \tiny QuIP\#  & \multicolumn{1}{l}{\tiny AQLM}  \\ \midrule
                                               & \multicolumn{1}{c|}{W2}                         & \multicolumn{1}{c|}{5.12}                         & \textbf{5.17} & \textbf{5.17} & \multicolumn{1}{c|}{5.22}                         & 5.19  & \multicolumn{1}{c|}{5.21}                         & \textbf{5.38} & 5.40          & \multicolumn{1}{c|}{5.60}                         & 5.41  & \multicolumn{1}{c|}{5.38}                         & 7.05          & \textbf{6.82} & \multicolumn{1}{c|}{8.22}                         & 6.19   & 6.14                      \\
\multirow{-2}{*}{2-7}                          & \multicolumn{1}{c|}{C4}                         & \multicolumn{1}{c|}{6.63}                         & \textbf{6.71} & \textbf{6.71} & \multicolumn{1}{c|}{6.79}                         & 6.75  & \multicolumn{1}{c|}{6.75}                         & \textbf{6.99} & 7.01          & \multicolumn{1}{c|}{7.34}                         & 7.04  & \multicolumn{1}{c|}{7.01}                         & 9.14          & \textbf{8.96} & \multicolumn{1}{c|}{11.0}                         & 8.16   & 8.09                      \\
\rowcolor[HTML]{D9D9D9} 
\cellcolor[HTML]{D9D9D9}                       & \multicolumn{1}{c|}{\cellcolor[HTML]{D9D9D9}W2} & \multicolumn{1}{c|}{\cellcolor[HTML]{D9D9D9}4.57} & \textbf{4.62} & \textbf{4.62} & \multicolumn{1}{c|}{\cellcolor[HTML]{D9D9D9}4.65} & 4.63  & \multicolumn{1}{c|}{\cellcolor[HTML]{D9D9D9}4.64} & \textbf{4.74} & \textbf{4.74} & \multicolumn{1}{c|}{\cellcolor[HTML]{D9D9D9}4.90} & 4.78  & \multicolumn{1}{c|}{\cellcolor[HTML]{D9D9D9}4.78} & 5.59          & \textbf{5.52} & \multicolumn{1}{c|}{\cellcolor[HTML]{D9D9D9}6.06} & 5.35   & 5.33                      \\
\rowcolor[HTML]{D9D9D9} 
\multirow{-2}{*}{\cellcolor[HTML]{D9D9D9}2-13} & \multicolumn{1}{c|}{\cellcolor[HTML]{D9D9D9}C4} & \multicolumn{1}{c|}{\cellcolor[HTML]{D9D9D9}6.05} & \textbf{6.10} & \textbf{6.10} & \multicolumn{1}{c|}{\cellcolor[HTML]{D9D9D9}6.15} & 6.13  & \multicolumn{1}{c|}{\cellcolor[HTML]{D9D9D9}6.14} & \textbf{6.28} & \textbf{6.28} & \multicolumn{1}{c|}{\cellcolor[HTML]{D9D9D9}6.50} & 6.35  & \multicolumn{1}{c|}{\cellcolor[HTML]{D9D9D9}6.33} & 7.46          & \textbf{7.39} & \multicolumn{1}{c|}{\cellcolor[HTML]{D9D9D9}8.07} & 7.20   & 7.19                      \\
                                               & \multicolumn{1}{c|}{W2}                         & \multicolumn{1}{c|}{3.12}                         & \textbf{3.16} & \textbf{3.16} & \multicolumn{1}{c|}{3.18}                         & 3.18  & \multicolumn{1}{c|}{3.19}                         & \textbf{3.27} & \textbf{3.27} & \multicolumn{1}{c|}{3.41}                         & 3.35  & \multicolumn{1}{c|}{3.36}                         & \textbf{3.87} & 3.90          & \multicolumn{1}{c|}{4.16}                         & 3.91   & 3.83                      \\
\multirow{-2}{*}{2-70}                         & \multicolumn{1}{c|}{C4}                         & \multicolumn{1}{c|}{4.97}                         & \textbf{5.00} & \textbf{5.00} & \multicolumn{1}{c|}{5.02}                         & 5.02  & \multicolumn{1}{c|}{5.03}                         & \textbf{5.09} & \textbf{5.09} & \multicolumn{1}{c|}{5.20}                         & 5.15  & \multicolumn{1}{c|}{5.17}                         & 5.70          & \textbf{5.69} & \multicolumn{1}{c|}{6.01}                         & 5.71   & 5.62                      \\ \bottomrule
\end{tabular}
\label{tab:llamacomp}
\end{table}

Here, we use 1MAD and 3INST with $L=16, V=1, T_x = T_y=16$.
Setting $T_x = T_y = 16$ enables using a $16\times16$ MMA tile per trellis sequence to perform matrix multiplication during inference.
$16 \times 16$ MMA tiles form the basis of many types of ``AI hardware,'' making fast decoding relatively simple \cite{9361255}.
We do not perform fine-tuning since the codes themselves are not tunable, but these codes are fully compatible with \qs-style fine-tuning (recall that \qs's codebook is also not tunable).
Table \ref{tab:llamacomp} shows that both 1MAD and 3INST significantly outperform \qss without fine-tuning (AQLM does not have numbers without fine-tuning).
Even at 4 bits, where all methods are close to lossless, \qts results in significant improvements.
Notably, the computed-code \qts variants \textit{without} fine-tuning outperforms both \qss and AQLM \textit{with} fine-tuning on almost all models and sizes, showing that fine-tuning is not a silver bullet.

\subsection{Hybrid Lookup-Computed Codes}
\begin{wraptable}{r}{6cm}
\caption{Batch size 1 decoding throughput on a RTX6000 Ada (960GB/s mem. BW).}
\small\sc
\centering
\tabcolsep=0.1cm
\begin{tabular}{@{}cccc@{}}
\toprule
Method & Bits & 2-7B Tok/s & 2-70B Tok/s \\ \midrule
FP16   & 16   & 55.9       & OOM         \\
AQLM   & 2    & 81.5       & 8.78        \\
QuIP\#  & 2    & 186        & 22.2        \\
\rowcolor[HTML]{D9D9D9} 
QTIP   & 2    & 188        & 23.5        \\
\rowcolor[HTML]{D9D9D9} 
QTIP   & 3    & 161        & 19.1        \\
\rowcolor[HTML]{D9D9D9} 
QTIP   & 4    & 140        & 16.3        \\ \bottomrule
\end{tabular}
\label{tab:genspeed}
\end{wraptable}

\label{sec:hybres}
\begin{table}[t]
\caption{Wikitext2 and C4 perplexity ($\downarrow$), \qts with the hybrid-computed code. \qts enables high-dimensional quantization and outperforms \sota vector quantization approaches.}
\small\sc
\tabcolsep=0.09cm
\centering
%\renewcommand{\arraystretch}{0.9}
\begin{tabular}{@{}cccccccccccccccc@{}}
\multicolumn{1}{l}{}              & \multicolumn{1}{l}{}                              & \multicolumn{8}{c}{ctx. 2048, \textcolor[HTML]{CC0000}{\textbf{X}} = GPTVQ, \textcolor{blue}{\textbf{Y}} = 0.13}                                                                                                                                                                      & \multicolumn{6}{c}{ctx. 4096, \textcolor[HTML]{CC0000}{\textbf{X}} = AQLM, \textcolor{blue}{\textbf{Y}} $\approx$ 0}                                                                                     \\ \midrule
\multicolumn{1}{l}{}              & \multicolumn{1}{l|}{}                             & \multicolumn{4}{c}{Wiktext2}                                                                               & \multicolumn{4}{c|}{C4}                                                                                    & \multicolumn{3}{c}{Wikitext2}                                                              & \multicolumn{3}{c}{C4}                        \\ \midrule
Method                            & \multicolumn{1}{c|}{Bits}                         & 1-7           & 1-13          & 1-30          & \multicolumn{1}{c|}{1-65}                                  & 1-7           & 1-13          & 1-30          & \multicolumn{1}{c|}{1-65}                                  & 2-7           & 2-13          & \multicolumn{1}{c|}{2-70}                                  & 2-7           & 2-13          & 2-70          \\ \midrule
FP16                              & \multicolumn{1}{c|}{16.0}                         & 5.68          & 5.09          & 4.10          & \multicolumn{1}{c|}{3.53}                                  & 7.04          & 6.61          & 5.98          & \multicolumn{1}{c|}{5.62}                                  & 5.12          & 4.57          & \multicolumn{1}{c|}{3.12}                                  & 6.63          & 6.05          & 4.97          \\
\rowcolor[HTML]{D9D9D9} 
{\color[HTML]{CC0000} \textbf{X}} & \multicolumn{1}{c|}{\cellcolor[HTML]{D9D9D9}4+\color{blue}{\textbf{Y}}}  & 5.94          & 5.20          & 4.18          & \multicolumn{1}{c|}{\cellcolor[HTML]{D9D9D9}3.64}          & --            & --            & --            & \multicolumn{1}{c|}{\cellcolor[HTML]{D9D9D9}--}            & 5.21          & 4.65          & \multicolumn{1}{c|}{\cellcolor[HTML]{D9D9D9}3.19}          & 6.75          & 6.14          & 5.03          \\
\rowcolor[HTML]{D9D9D9} 
QuIP\#                            & \multicolumn{1}{c|}{\cellcolor[HTML]{D9D9D9}4.00} & 5.76          & 5.17          & 4.18          & \multicolumn{1}{c|}{\cellcolor[HTML]{D9D9D9}3.60}          & 7.18          & 6.67          & 6.03          & \multicolumn{1}{c|}{\cellcolor[HTML]{D9D9D9}5.66}          & 5.19          & 4.63          & \multicolumn{1}{c|}{\cellcolor[HTML]{D9D9D9}3.18}          & 6.75          & 6.13          & 5.02          \\
\rowcolor[HTML]{D9D9D9} 
QTIP                              & \multicolumn{1}{c|}{\cellcolor[HTML]{D9D9D9}4.00} & \textbf{5.72} & \textbf{5.15} & \textbf{4.15} & \multicolumn{1}{c|}{\cellcolor[HTML]{D9D9D9}\textbf{3.58}} & \textbf{7.13} & \textbf{6.65} & \textbf{6.01} & \multicolumn{1}{c|}{\cellcolor[HTML]{D9D9D9}\textbf{5.64}} & \textbf{5.17} & \textbf{4.61} & \multicolumn{1}{c|}{\cellcolor[HTML]{D9D9D9}\textbf{3.16}} & \textbf{6.69} & \textbf{6.09} & \textbf{5.00} \\
{\color[HTML]{CC0000} \textbf{X}} & \multicolumn{1}{c|}{3+\color{blue}{\textbf{Y}}}                          & 6.32          & 5.31          & 4.38          & \multicolumn{1}{c|}{3.79}                                  & --            & --            & --            & \multicolumn{1}{c|}{--}                                    & 5.38          & 4.78          & \multicolumn{1}{c|}{3.36}                                  & 7.01          & 6.33          & 5.17          \\
QuIP\#                            & \multicolumn{1}{c|}{3.00}                         & 5.98          & 5.31          & 4.36          & \multicolumn{1}{c|}{3.70}                                  & 7.39          & 6.83          & 6.17          & \multicolumn{1}{c|}{5.77}                                  & 5.41          & 4.78          & \multicolumn{1}{c|}{3.35}                                  & 7.04          & 6.35          & 5.15          \\
QTIP                              & \multicolumn{1}{c|}{3.00}                         & \textbf{5.85} & \textbf{5.24} & \textbf{4.26} & \multicolumn{1}{c|}{\textbf{3.68}}                         & \textbf{7.26} & \textbf{6.74} & \textbf{6.09} & \multicolumn{1}{c|}{\textbf{5.71}}                         & \textbf{5.28} & \textbf{4.69} & \multicolumn{1}{c|}{\textbf{3.26}}                         & \textbf{6.87} & \textbf{6.22} & \textbf{5.08} \\
\rowcolor[HTML]{D9D9D9} 
{\color[HTML]{CC0000} \textbf{X}} & \multicolumn{1}{c|}{\cellcolor[HTML]{D9D9D9}2+\color{blue}{\textbf{Y}}}  & 9.64          & 6.58          & 5.63          & \multicolumn{1}{c|}{\cellcolor[HTML]{D9D9D9}4.91}          & --            & --            & --            & \multicolumn{1}{c|}{\cellcolor[HTML]{D9D9D9}--}            & 6.14          & 5.33          & \multicolumn{1}{c|}{\cellcolor[HTML]{D9D9D9}3.83}          & 8.09          & 7.19          & 5.62          \\
\rowcolor[HTML]{D9D9D9} 
QuIP\#                            & \multicolumn{1}{c|}{\cellcolor[HTML]{D9D9D9}2.00} & 6.86          & 5.97          & 5.02          & \multicolumn{1}{c|}{\cellcolor[HTML]{D9D9D9}4.36}          & 8.36          & 7.48          & 6.71          & \multicolumn{1}{c|}{\cellcolor[HTML]{D9D9D9}6.19}          & 6.19          & 5.35          & \multicolumn{1}{c|}{\cellcolor[HTML]{D9D9D9}3.91}          & 8.16          & 7.20          & 5.71          \\
\rowcolor[HTML]{D9D9D9} 
QTIP                              & \multicolumn{1}{c|}{\cellcolor[HTML]{D9D9D9}2.00} & \textbf{6.52} & \textbf{5.80} & \textbf{4.83} & \multicolumn{1}{c|}{\cellcolor[HTML]{D9D9D9}\textbf{4.21}} & \textbf{7.99} & \textbf{7.31} & \textbf{6.56} & \multicolumn{1}{c|}{\cellcolor[HTML]{D9D9D9}\textbf{6.08}} & \textbf{5.86} & \textbf{5.11} & \multicolumn{1}{c|}{\cellcolor[HTML]{D9D9D9}\textbf{3.70}} & \textbf{7.73} & \textbf{6.85} & \textbf{5.48} \\ \bottomrule
\end{tabular}
\label{tab:hybft}
\end{table}



% \begin{tabular}{@{}cccccccccccc@{}}
%                        &                                                   &                                                   & \multicolumn{3}{c}{$\sim$4 Bit}                                                                                 & \multicolumn{3}{c}{$\sim$3 Bit}                                                                                 & \multicolumn{3}{c}{$\sim$2 Bit}                                                            \\ \midrule
%                        & \multicolumn{1}{c|}{}                             & \multicolumn{1}{c|}{FP16}                         & QTIP                         & QuIP\#                       & \multicolumn{1}{c|}{AQLM}                         & QTIP                         & QuIP\#                       & \multicolumn{1}{c|}{AQLM}                         & QTIP                         & QuIP\#                       & AQLM                         \\ \midrule
%                        & \multicolumn{1}{c|}{\cellcolor[HTML]{D9D9D9}Bits} & \multicolumn{1}{c|}{\cellcolor[HTML]{D9D9D9}16.0} & \cellcolor[HTML]{D9D9D9}4.00 & \cellcolor[HTML]{D9D9D9}4.00 & \multicolumn{1}{c|}{\cellcolor[HTML]{D9D9D9}4.04} & \cellcolor[HTML]{D9D9D9}3.00 & \cellcolor[HTML]{D9D9D9}3.00 & \multicolumn{1}{c|}{\cellcolor[HTML]{D9D9D9}3.04} & \cellcolor[HTML]{D9D9D9}2.00 & \cellcolor[HTML]{D9D9D9}2.00 & \cellcolor[HTML]{D9D9D9}2.02 \\
%                        & \multicolumn{1}{c|}{W2}                           & \multicolumn{1}{c|}{5.12}                         & \textbf{5.17}                & 5.19                         & \multicolumn{1}{c|}{5.21}                         & \textbf{5.29}                & 5.41                         & \multicolumn{1}{c|}{5.46}                         & \textbf{5.91}                & 6.19                         & 6.64                         \\
% \multirow{-3}{*}{2-7}  & \multicolumn{1}{c|}{C4}                           & \multicolumn{1}{c|}{6.63}                         & \textbf{6.69}                & 6.75                         & \multicolumn{1}{c|}{6.75}                         & \textbf{6.88}                & 7.04                         & \multicolumn{1}{c|}{7.08}                         & \textbf{7.76}                & 8.16                         & 8.56                         \\ \midrule
%                        & \multicolumn{1}{c|}{\cellcolor[HTML]{D9D9D9}Bits} & \multicolumn{1}{c|}{\cellcolor[HTML]{D9D9D9}16.0} & \cellcolor[HTML]{D9D9D9}4.00 & \cellcolor[HTML]{D9D9D9}4.00 & \multicolumn{1}{c|}{\cellcolor[HTML]{D9D9D9}3.94} & \cellcolor[HTML]{D9D9D9}3.00 & \cellcolor[HTML]{D9D9D9}3.00 & \multicolumn{1}{c|}{\cellcolor[HTML]{D9D9D9}3.03} & \cellcolor[HTML]{D9D9D9}2.00 & \cellcolor[HTML]{D9D9D9}2.00 & \cellcolor[HTML]{D9D9D9}1.97 \\
%                        & \multicolumn{1}{c|}{W2}                           & \multicolumn{1}{c|}{4.57}                         & \textbf{4.61}                & 4.63                         & \multicolumn{1}{c|}{4.65}                         & \textbf{4.71}                & 4.78                         & \multicolumn{1}{c|}{4.82}                         & \textbf{5.26}                & 5.35                         & 5.65                         \\
% \multirow{-3}{*}{2-13} & \multicolumn{1}{c|}{C4}                           & \multicolumn{1}{c|}{6.05}                         & \textbf{6.09}                & 6.13                         & \multicolumn{1}{c|}{6.14}                         & \textbf{6.23}                & 6.35                         & \multicolumn{1}{c|}{6.37}                         & \textbf{6.99}                & 7.20                         & 7.51                         \\ \midrule
%                        & \multicolumn{1}{c|}{\cellcolor[HTML]{D9D9D9}Bits} & \multicolumn{1}{c|}{\cellcolor[HTML]{D9D9D9}16.0} & \cellcolor[HTML]{D9D9D9}4.00 & \cellcolor[HTML]{D9D9D9}4.00 & \multicolumn{1}{c|}{\cellcolor[HTML]{D9D9D9}4.14} & \cellcolor[HTML]{D9D9D9}3.00 & \cellcolor[HTML]{D9D9D9}3.00 & \multicolumn{1}{c|}{\cellcolor[HTML]{D9D9D9}3.01} & \cellcolor[HTML]{D9D9D9}2.00 & \cellcolor[HTML]{D9D9D9}2.00 & \cellcolor[HTML]{D9D9D9}2.07 \\
%                        & \multicolumn{1}{c|}{W2}                           & \multicolumn{1}{c|}{3.12}                         & \textbf{3.16}                & 3.18                         & \multicolumn{1}{c|}{3.19}                         & \textbf{3.26}                & 3.35                         & \multicolumn{1}{c|}{3.36}                         & \textbf{3.78}                & 3.91                         & 3.94                         \\
% \multirow{-3}{*}{2-70} & \multicolumn{1}{c|}{C4}                           & \multicolumn{1}{c|}{4.97}                         & \textbf{5.00}                & 5.02                         & \multicolumn{1}{c|}{5.03}                         & \textbf{5.08}                & 5.15                         & \multicolumn{1}{c|}{5.17}                         & \textbf{5.56}                & 5.71                         & 5.72                         \\ \bottomrule
% \end{tabular}

% \begin{tabular}{@{}ccclccclccclccc@{}}
% \multicolumn{1}{l}{}   & \multicolumn{1}{l}{}                              & \multicolumn{1}{l}{}                               & \multicolumn{4}{c}{$\sim$4 Bit}                                                                                                                                    & \multicolumn{4}{c}{$\sim$3 Bit}                                                                                                                                    & \multicolumn{4}{c}{$\sim$2 Bit}                                                                                                               \\ \midrule
% \multicolumn{1}{l}{}   & \multicolumn{1}{l|}{}                             & \multicolumn{1}{c|}{\rotatebox{45}{FP16}}                          & \multicolumn{1}{c}{\rotatebox{45}{QTIP}}                         & \rotatebox{45}{QuIP\#}                       & \rotatebox{45}{gptvq}                        & \multicolumn{1}{c|}{\rotatebox{45}{omniq}}                        & \multicolumn{1}{c}{\rotatebox{45}{QTIP}}                         & \rotatebox{45}{QuIP\#}                       & \rotatebox{45}{gptvq}                        & \multicolumn{1}{c|}{\rotatebox{45}{omniq}}                        & \multicolumn{1}{c}{\rotatebox{45}{QTIP}}                         & \rotatebox{45}{QuIP\#}                       & \rotatebox{45}{gptvq}                        & \rotatebox{45}{omniq}                        \\ \midrule
% \multicolumn{1}{l}{}   & \multicolumn{1}{c|}{\cellcolor[HTML]{D9D9D9}Bits} & \multicolumn{1}{c|}{\cellcolor[HTML]{D9D9D9}16.0} & \multicolumn{1}{c}{\cellcolor[HTML]{D9D9D9}4.00} & \cellcolor[HTML]{D9D9D9}4.00 & \cellcolor[HTML]{D9D9D9}4.13 & \multicolumn{1}{c|}{\cellcolor[HTML]{D9D9D9}4.13} & \multicolumn{1}{c}{\cellcolor[HTML]{D9D9D9}3.00} & \cellcolor[HTML]{D9D9D9}3.00 & \cellcolor[HTML]{D9D9D9}3.13 & \multicolumn{1}{c|}{\cellcolor[HTML]{D9D9D9}3.13} & \multicolumn{1}{c}{\cellcolor[HTML]{D9D9D9}2.00} & \cellcolor[HTML]{D9D9D9}2.00 & \cellcolor[HTML]{D9D9D9}2.25 & \cellcolor[HTML]{D9D9D9}2.25 \\ \midrule
%                        & \multicolumn{1}{c|}{W2}                           & \multicolumn{1}{c|}{5.68}                          &                                                  & 5.76                         & 5.94                         & \multicolumn{1}{c|}{5.77}                         & \multicolumn{1}{c}{5.87}                         & 5.98                         & 6.32                         & \multicolumn{1}{c|}{6.15}                         & \multicolumn{1}{c}{6.65}                         & 6.86                         & 8.76                         & 8.9                          \\
% \multirow{-2}{*}{1-7}  & \multicolumn{1}{c|}{C4}                           & \multicolumn{1}{c|}{7.04}                          &                                                  & 7.18                         & --                           & \multicolumn{1}{c|}{7.21}                         & \multicolumn{1}{c}{7.26}                         & 7.39                         & --                           & \multicolumn{1}{c|}{7.75}                         & \multicolumn{1}{c}{8.11}                         & 8.36                         & --                           & 11.78                        \\ \midrule
%                        & \multicolumn{1}{c|}{W2}                           & \multicolumn{1}{c|}{5.09}                          &                                                  & 5.17                         & 5.2                          & \multicolumn{1}{c|}{5.17}                         &                                                  & 5.31                         & 5.31                         & \multicolumn{1}{c|}{5.44}                         &                                                  & 5.97                         & 6.33                         & 7.34                         \\
% \multirow{-2}{*}{1-13} & \multicolumn{1}{c|}{C4}                           & \multicolumn{1}{c|}{6.61}                          &                                                  & 6.67                         & --                           & \multicolumn{1}{c|}{6.69}                         &                                                  & 6.83                         & --                           & \multicolumn{1}{c|}{7.05}                         &                                                  & 7.48                         & --                           & 9.75                         \\ \midrule
%                        & \multicolumn{1}{c|}{W2}                           & \multicolumn{1}{c|}{4.10}                          &                                                  & 4.18                         & 4.18                         & \multicolumn{1}{c|}{4.19}                         &                                                  & 4.36                         & 4.38                         & \multicolumn{1}{c|}{4.56}                         &                                                  & 5.02                         & 5.42                         & 6.59                         \\
% \multirow{-2}{*}{1-30} & \multicolumn{1}{c|}{C4}                           & \multicolumn{1}{c|}{5.98}                          &                                                  & 6.03                         & --                           & \multicolumn{1}{c|}{6.06}                         &                                                  & 6.17                         & --                           & \multicolumn{1}{c|}{6.37}                         &                                                  & 6.71                         & --                           & 8.65                         \\ \midrule
%                        & \multicolumn{1}{c|}{W2}                           & \multicolumn{1}{c|}{3.53}                          &                                                  & 3.60                         & 3.64                         & \multicolumn{1}{c|}{3.62}                         &                                                  & 3.70                         & 3.79                         & \multicolumn{1}{c|}{3.94}                         &                                                  & 4.36                         & 4.74                         & 5.65                         \\
% \multirow{-2}{*}{1-65} & \multicolumn{1}{c|}{C4}                           & \multicolumn{1}{c|}{5.62}                          &                                                  & 5.66                         & --                           & \multicolumn{1}{c|}{5.68}                         &                                                  & 5.77                         & --                           & \multicolumn{1}{c|}{5.93}                         &                                                  & 6.19                         & --                           & 7.6                          \\ \bottomrule
% \end{tabular}

% $\hspace{1.5cm}\xleftarrow{\hspace*{4cm}\mbox{generally better}\hspace*{4cm}}$

\begin{table}[t]
\caption{Zeroshot accuracy ($\uparrow$), \qts with the hybrid-computed code.}
\small\sc
\tabcolsep=0.04cm
%\renewcommand{\arraystretch}{0.9}
\centering
\begin{tabular}{@{}cccccccccccccccc@{}}
                                                    & \multicolumn{5}{c}{2-70}                                                                                          & \multicolumn{5}{c}{2-13}                                                                                          & \multicolumn{5}{c}{2-7}                                              \\ \midrule
\multicolumn{1}{c|}{Mthd.}                          & Bits & ArcC          & ArcE          & PiQA          & \multicolumn{1}{c|}{Wino}                                  & Bits & ArcC          & ArcE          & PiQA          & \multicolumn{1}{c|}{Wino}                                  & Bits & ArcC          & ArcE          & PiQA          & Wino          \\ \midrule
\multicolumn{1}{c|}{FP16}                           & 16   & 51.1          & 77.7          & 81.1          & \multicolumn{1}{c|}{77.0}                                  & 16   & 45.6          & 73.3          & 73.5          & \multicolumn{1}{c|}{69.6}                                  & 16   & 40.0          & 69.3          & 78.5          & 67.3          \\
\rowcolor[HTML]{D9D9D9} 
\multicolumn{1}{c|}{\cellcolor[HTML]{D9D9D9}AQLM}   & 4.14 & 50.7          & 77.3          & 81.5          & \multicolumn{1}{c|}{\cellcolor[HTML]{D9D9D9}76.5}          & 3.94 & \textbf{44.8} & 73.3          & 78.4          & \multicolumn{1}{c|}{\cellcolor[HTML]{D9D9D9}\textbf{69.9}} & 4.04 & 41.0          & 70.2          & 78.2          & 67.3          \\
\rowcolor[HTML]{D9D9D9} 
\multicolumn{1}{c|}{\cellcolor[HTML]{D9D9D9}QuIP\#} & 4    & 50.5          & 77.7          & 81.4          & \multicolumn{1}{c|}{\cellcolor[HTML]{D9D9D9}\textbf{77.3}} & 4    & 43.6          & 71.3          & 78.7          & \multicolumn{1}{c|}{\cellcolor[HTML]{D9D9D9}69.6}          & 4    & 40.4          & 68.6          & \textbf{78.5} & \textbf{67.4} \\
\rowcolor[HTML]{D9D9D9} 
\multicolumn{1}{c|}{\cellcolor[HTML]{D9D9D9}QTIP}   & 4    & 50.0          & \textbf{77.8} & 81.3          & \multicolumn{1}{c|}{\cellcolor[HTML]{D9D9D9}76.9}          & 4    & 44.8          & 73.6          & \textbf{78.9} & \multicolumn{1}{c|}{\cellcolor[HTML]{D9D9D9}69.9}          & 4    & 40.0          & \textbf{68.9} & 78.4          & 67.1          \\
\multicolumn{1}{c|}{AQLM}                           & 3.01 & 50.3          & 78.0          & 80.7          & \multicolumn{1}{c|}{75.3}                                  & 3.03 & 42.8          & 72.9          & 78.5          & \multicolumn{1}{c|}{68.8}                                  & 3.04 & 38.5          & 66.8          & 77.3          & 65.4          \\
\multicolumn{1}{c|}{QuIP\#}                         & 3    & \textbf{50.9} & 77.6          & \textbf{81.4} & \multicolumn{1}{c|}{76.1}                                  & 3    & \textbf{44.0} & 72.5          & 78.4          & \multicolumn{1}{c|}{69.1}                                  & 3    & \textbf{39.2} & \textbf{68.4} & 77.3          & 66.5          \\
\multicolumn{1}{c|}{QTIP}                           & 3    & 50.3          & \textbf{78.2} & 80.6          & \multicolumn{1}{c|}{77.0}                                  & 3    & \textbf{44.0} & 72.8          & 78.0          & \multicolumn{1}{c|}{\textbf{69.5}}                         & 3    & 38.9          & 68.1          & \textbf{78.1} & \textbf{66.9} \\
\rowcolor[HTML]{D9D9D9} 
\multicolumn{1}{c|}{\cellcolor[HTML]{D9D9D9}AQLM}   & 2.07 & 47.9          & \textbf{77.7} & 80.4          & \multicolumn{1}{c|}{\cellcolor[HTML]{D9D9D9}75.9}          & 1.97 & 38.8          & 69.3          & 75.9          & \multicolumn{1}{c|}{\cellcolor[HTML]{D9D9D9}68.8}          & 2.02 & 32.8          & 63.7          & 74.8          & 65.7          \\
\rowcolor[HTML]{D9D9D9} 
\multicolumn{1}{c|}{\cellcolor[HTML]{D9D9D9}QuIP\#} & 2    & 47.6          & 77.1          & 79.5          & \multicolumn{1}{c|}{\cellcolor[HTML]{D9D9D9}74.6}          & 2    & 39.6          & 69.0          & \textbf{77.3} & \multicolumn{1}{c|}{\cellcolor[HTML]{D9D9D9}67.4}          & 2    & 35.2          & 65.3          & 75.4          & \textbf{64.9} \\
\rowcolor[HTML]{D9D9D9} 
\multicolumn{1}{c|}{\cellcolor[HTML]{D9D9D9}QTIP}   & 2    & \textbf{48.0} & 76.3          & 80.2          & \multicolumn{1}{c|}{\cellcolor[HTML]{D9D9D9}75.1}          & 2    & \textbf{41.4} & \textbf{70.8} & \textbf{77.3} & \multicolumn{1}{c|}{\cellcolor[HTML]{D9D9D9}\textbf{67.6}} & 2    & \textbf{35.7} & \textbf{65.6} & \textbf{75.9} & 64.7          \\ \bottomrule
\end{tabular}
\label{tab:hybftzs}
\end{table}
\begin{table}[t]
\caption{\qts vs. \qs, Llama 3 (ctx. 8192 for perplexity). The 70B results quantize layer 0 v, which catastrophically degrades zeroshot performance without special prefix tokens (e.g. \texttt{\textbackslash n}). This can be remedied by prepending the prompt with \texttt{\textbackslash n}, applying the chat template as in Table 8, or not quantizing 0 v (see \cite{yaqa}). Regardless, the lower distortion of TCQ in \qts improves over the low-dimensional VQ in \qs.}
\small\sc
\tabcolsep=0.09cm
%\renewcommand{\arraystretch}{0.9}
\centering
\begin{tabular}{@{}cccccccccccccccc@{}}
\multicolumn{1}{l}{} & \multicolumn{1}{l}{}                              & \multicolumn{2}{c}{3-70 ppl ($\downarrow$)}                                & \multicolumn{5}{c}{3-70 zeroshot acc ($\uparrow$)}                                                                        & \multicolumn{2}{c}{3-8 ppl ($\downarrow$)}                                 & \multicolumn{5}{c}{3-8 zeroshot acc ($\uparrow$)}                            \\ \midrule
Mthd.               & \multicolumn{1}{c|}{Bits}                         & W2            & \multicolumn{1}{c|}{C4}                                    & \scriptsize ArcC          & \scriptsize ArcE          & \scriptsize BoolQ         & \scriptsize PiQA          & \multicolumn{1}{c|}{\scriptsize Wino}                                  & W2            & \multicolumn{1}{c|}{C4}                                    & \scriptsize ArcC          & \scriptsize ArcE          & \scriptsize BoolQ         & \scriptsize PiQA          & \scriptsize Wino          \\ \midrule
BF16                 & \multicolumn{1}{c|}{16.0}                         & 2.59          & \multicolumn{1}{c|}{5.78}                                  & 60.5          & 86.9          & 85.3          & 82.4          & \multicolumn{1}{c|}{80.3}                                  & 5.54          & \multicolumn{1}{c|}{7.10}                                  & 50.2          & 80.1          & 81.0          & 79.7          & 72.9          \\
\rowcolor[HTML]{D9D9D9} 
QuIP\#               & \multicolumn{1}{c|}{\cellcolor[HTML]{D9D9D9}4.00} & 2.99          & \multicolumn{1}{c|}{\cellcolor[HTML]{D9D9D9}5.96}          & 35.0          & 67.3          & 84.7          & 71.9          & \multicolumn{1}{c|}{\cellcolor[HTML]{D9D9D9}76.7}          & 5.81          & \multicolumn{1}{c|}{\cellcolor[HTML]{D9D9D9}7.32}          & 50.2          & \textbf{79.7} & \textbf{81.3} & \textbf{79.7} & 73.1          \\
\rowcolor[HTML]{D9D9D9} 
QTIP                 & \multicolumn{1}{c|}{\cellcolor[HTML]{D9D9D9}4.00} & \textbf{2.75} & \multicolumn{1}{c|}{\cellcolor[HTML]{D9D9D9}5.83}          & \textbf{56.1} & \textbf{83.9} & \textbf{85.8} & \textbf{81.3} & \multicolumn{1}{c|}{\cellcolor[HTML]{D9D9D9}\textbf{80.6}} & \textbf{5.67} & \multicolumn{1}{c|}{\cellcolor[HTML]{D9D9D9}\textbf{7.20}} & 50.2          & 79.6          & 79.5          & 79.4          & \textbf{73.4} \\
QuIP\#               & \multicolumn{1}{c|}{3.00}                         & 3.59          & \multicolumn{1}{c|}{6.18}                                  & 31.1          & 36.6          & \textbf{85.7} & 58.8          & \multicolumn{1}{c|}{76.4}                                  & 6.27          & \multicolumn{1}{c|}{7.71}                                  & 46.4          & 77.4          & 79.9          & 77.9          & 72.9          \\
QTIP                 & \multicolumn{1}{c|}{3.00}                         & \textbf{3.18} & \multicolumn{1}{c|}{5.98}                                  & \textbf{48.6} & \textbf{77.8} & 85.0          & \textbf{77.8} & \multicolumn{1}{c|}{\textbf{79.7}}                         & \textbf{6.01} & \multicolumn{1}{c|}{\textbf{7.48}}                         & \textbf{49.2} & \textbf{79.3} & \textbf{80.0} & \textbf{79.2} & \textbf{74.5} \\
\rowcolor[HTML]{D9D9D9} 
QuIP\#               & \multicolumn{1}{c|}{\cellcolor[HTML]{D9D9D9}2.00} & 5.77          & \multicolumn{1}{c|}{\cellcolor[HTML]{D9D9D9}7.46}          & 18.3          & 32.2 & 82.1          & 54.7          & \multicolumn{1}{c|}{\cellcolor[HTML]{D9D9D9}68.9}          & 7.84          & \multicolumn{1}{c|}{\cellcolor[HTML]{D9D9D9}9.06}          & 39.2          & 72.9          & 76.6          & 75.6          & 68.2          \\
\rowcolor[HTML]{D9D9D9} 
QTIP                 & \multicolumn{1}{c|}{\cellcolor[HTML]{D9D9D9}2.00} & \textbf{4.97} & \multicolumn{1}{c|}{\cellcolor[HTML]{D9D9D9}\textbf{6.80}} & \textbf{28.0} & \textbf{35.2}          & \textbf{83.6} & \textbf{57.1} & \multicolumn{1}{c|}{\cellcolor[HTML]{D9D9D9}\textbf{72.6}} & \textbf{7.33} & \multicolumn{1}{c|}{\cellcolor[HTML]{D9D9D9}\textbf{8.62}} & \textbf{44.2} & \textbf{75.2} & \textbf{76.7} & \textbf{77.6} & \textbf{70.7} \\ \bottomrule
\end{tabular}
\label{tab:llama3}
\end{table}
\begin{table}[t]
\caption{Llama 3.1 instruct-tuned model results (ctx. 8192 for perplexity).  QTIP performs well at all model sizes and generally outperforms PV-Tuning, a recent quantization method that focuses on fine-tuning algorithms. The zeroshot results in this table use LM Eval 0.4.4 and the ``standard'' versions of each task instead of the Meta versions in \cite{llama3}.}
\label{tab:llama31}
\small\sc\centering\tabcolsep=0.1cm
\begin{tabular}{@{}cccccccc@{}}
                                                                             &                                                        &                                                   & Ppl. ($\downarrow$)                                 & \multicolumn{4}{c}{Zeroshot ($\uparrow$)} \\ \midrule
                                                                             & \multicolumn{1}{c|}{}                                  & \multicolumn{1}{c|}{Bits}                         & \multicolumn{1}{c|}{W2}                           & ArcC      & ArcE     & Hswag     & PiQA     \\ \midrule
\multicolumn{1}{c|}{}                                                        & \multicolumn{1}{c|}{Meta "FP8"}                        & \multicolumn{1}{c|}{16 Attn. / 8 MLP}             & \multicolumn{1}{c|}{1.70}                         & 61.6      & 81.4     & 67.1      & 83.8     \\
\multicolumn{1}{c|}{}                                                        & \multicolumn{1}{c|}{QTIP}                              & \multicolumn{1}{c|}{4}                            & \multicolumn{1}{c|}{1.79}                         & 61.3      & 80.9     & 66.7      & 84.2     \\
\multicolumn{1}{c|}{}                                                        & \multicolumn{1}{c|}{QTIP}                              & \multicolumn{1}{c|}{3}                            & \multicolumn{1}{c|}{2.05}                         & 61.5      & 81.4     & 66.8      & 83.5     \\
\multicolumn{1}{c|}{\multirow{-4}{*}{3.1 405B Inst.}}                        & \multicolumn{1}{c|}{QTIP}                              & \multicolumn{1}{c|}{2}                            & \multicolumn{1}{c|}{3.29}                         & 60.7      & 81.1     & 65.4      & 82.2     \\
\rowcolor[HTML]{D9D9D9} 
\multicolumn{1}{c|}{\cellcolor[HTML]{D9D9D9}}                                & \multicolumn{1}{c|}{\cellcolor[HTML]{D9D9D9}BF16}      & \multicolumn{1}{c|}{\cellcolor[HTML]{D9D9D9}16}   & \multicolumn{1}{c|}{\cellcolor[HTML]{D9D9D9}3.52} & 56.7      & 75.6     & 61.5      & 82.8     \\
\rowcolor[HTML]{D9D9D9} 
\multicolumn{1}{c|}{\cellcolor[HTML]{D9D9D9}}                                & \multicolumn{1}{c|}{\cellcolor[HTML]{D9D9D9}QTIP}      & \multicolumn{1}{c|}{\cellcolor[HTML]{D9D9D9}4}    & \multicolumn{1}{c|}{\cellcolor[HTML]{D9D9D9}3.73} & 56.3      & 75.8     & 61.4      & 83.0     \\
\rowcolor[HTML]{D9D9D9} 
\multicolumn{1}{c|}{\cellcolor[HTML]{D9D9D9}}                                & \multicolumn{1}{c|}{\cellcolor[HTML]{D9D9D9}QTIP}      & \multicolumn{1}{c|}{\cellcolor[HTML]{D9D9D9}3}    & \multicolumn{1}{c|}{\cellcolor[HTML]{D9D9D9}4.12} & 55.1      & 75.1     & 60.8      & 82.6     \\
\rowcolor[HTML]{D9D9D9} 
\multicolumn{1}{c|}{\cellcolor[HTML]{D9D9D9}}                                & \multicolumn{1}{c|}{\cellcolor[HTML]{D9D9D9}QTIP}      & \multicolumn{1}{c|}{\cellcolor[HTML]{D9D9D9}2}    & \multicolumn{1}{c|}{\cellcolor[HTML]{D9D9D9}5.08} & 54.4      & 72.6     & 59.4      & 82.5     \\
\rowcolor[HTML]{D9D9D9} 
\multicolumn{1}{c|}{\multirow{-5}{*}{\cellcolor[HTML]{D9D9D9}3.1 70B Inst.}} & \multicolumn{1}{c|}{\cellcolor[HTML]{D9D9D9}PV-Tuning} & \multicolumn{1}{c|}{\cellcolor[HTML]{D9D9D9}2.01} & \multicolumn{1}{c|}{\cellcolor[HTML]{D9D9D9}5.70} & 52.7      & 72.2     & 60.2      & 82.6     \\
\multicolumn{1}{c|}{}                                                        & \multicolumn{1}{c|}{BF16}                              & \multicolumn{1}{c|}{16}                           & \multicolumn{1}{c|}{6.50}                         & 51.6      & 77.8     & 57.7      & 80.0     \\
\multicolumn{1}{c|}{}                                                        & \multicolumn{1}{c|}{QTIP}                              & \multicolumn{1}{c|}{4}                            & \multicolumn{1}{c|}{6.61}                         & 50.7      & 78.0     & 57.5      & 80.1     \\
\multicolumn{1}{c|}{}                                                        & \multicolumn{1}{c|}{QTIP}                              & \multicolumn{1}{c|}{3}                            & \multicolumn{1}{c|}{6.80}                         & 50.4      & 77.7     & 56.9      & 79.3     \\
\multicolumn{1}{c|}{}                                                        & \multicolumn{1}{c|}{QTIP}                              & \multicolumn{1}{c|}{2}                            & \multicolumn{1}{c|}{7.82}                         & 45.1      & 75.6     & 54.5      & 79.0     \\
\multicolumn{1}{c|}{\multirow{-5}{*}{3.1 8B Inst.}}                          & \multicolumn{1}{c|}{PV-Tuning}                         & \multicolumn{1}{c|}{2.07}                         & \multicolumn{1}{c|}{8.45}                         & 46.2      & 75.4     & 54.4      & 78.7     \\ \bottomrule
\end{tabular}
\end{table}
\begin{table}[ht!]
\caption{Llama 3.2 instruct-tuned results when quantizing to 4 bits (ctx. 8192 for perplexity). Even on extremely small models, QTIP is still able to achieve meaningful compression without sacrificing quality. This table uses the same LM Eval setup as Table \ref{tab:llama31}.}
\label{tab:llama32}
\small\sc\centering
\begin{tabular}{@{}cccccccc@{}}
                                                                  &                                                   &                                                   & Ppl ($\downarrow$)                                 & \multicolumn{4}{c}{Zeroshot ($\uparrow$)} \\ \midrule
                                                                  & \multicolumn{1}{c|}{}                             & \multicolumn{1}{c|}{Size (GB)}                    & \multicolumn{1}{c|}{W2}                            & ArcC     & ArcE     & Hswag     & PiQA    \\ \midrule
\multicolumn{1}{c|}{}                                             & \multicolumn{1}{c|}{BF16}                         & \multicolumn{1}{c|}{6}                            & \multicolumn{1}{c|}{9.58}                          & 43.3     & 74.3     & 52.2      & 75.7    \\
\multicolumn{1}{c|}{\multirow{-2}{*}{3B}}                         & \multicolumn{1}{c|}{QTIP}                         & \multicolumn{1}{c|}{2.1}                          & \multicolumn{1}{c|}{9.77}                          & 43.5     & 74.3     & 51.9      & 75.1    \\
\rowcolor[HTML]{D9D9D9} 
\multicolumn{1}{c|}{\cellcolor[HTML]{D9D9D9}}                     & \multicolumn{1}{c|}{\cellcolor[HTML]{D9D9D9}BF16} & \multicolumn{1}{c|}{\cellcolor[HTML]{D9D9D9}2.4}  & \multicolumn{1}{c|}{\cellcolor[HTML]{D9D9D9}11.57} & 36.0     & 68.5     & 45.2      & 74.2    \\
\rowcolor[HTML]{D9D9D9} 
\multicolumn{1}{c|}{\multirow{-2}{*}{\cellcolor[HTML]{D9D9D9}1B}} & \multicolumn{1}{c|}{\cellcolor[HTML]{D9D9D9}QTIP} & \multicolumn{1}{c|}{\cellcolor[HTML]{D9D9D9}0.97} & \multicolumn{1}{c|}{\cellcolor[HTML]{D9D9D9}11.93} & 34.8     & 68.4     & 44.5      & 73.3    \\ \bottomrule
\end{tabular}
\end{table}

Here, we use the hybrid lookup-computed code with $L=16, V=2, T_x = T_y=16, Q=9$.
Setting $Q=9$ gives a 2KiB codebook, which fits in L1 cache \textit{even after} duplication for bank conflicts (32$\times$) on modern GPUs.
This codebook is differentiable, so we can fine-tune it: to evaluate this, we fine-tune using \qs's methodology, tuning both the codebook entries and the as-yet-unquantized weights in a blockwise fashion.
Table \ref{tab:hybft} shows the perplexity of quantized Llama 1 and 2 models.
In all cases, \qts outperforms the other vector quantization-based methods.
Even at 3 and 4 bits, where \qss and AQLM are close to lossless, \qts roughly \textit{halves} the perplexity gap. 
These results also show the importance of dimensionality.
Note that the 3- and 4-bit Llama 2 70B numbers here match those in \ref{tab:llamacomp}.
Since Table \ref{tab:llamacomp} uses a pure-computed code \textit{without fine-tuning}, fine-tuning has no effect in these regimes and the improvement over \qss is purely from dimensionality.

Table \ref{tab:hybftzs} shows zeroshot results computed with LM Eval, which are slightly random; \qts generally matches or exceeds \qss and AQLM on these tasks \cite{lmeval}.
Table \ref{tab:llama3} contains results on Llama 3.
In our experiments, we observed that quantizing layer 0 v of all Llama 3 70B variants (including 3.1, 3.3, and all instruct models) resulted in catastrophic collapse of zeroshot performance. 
Since the focus of this work is on \textit{what to round with} and not \textit{how to round}, Table \ref{tab:llama3} includes results \textit{with} quantizing 0 v. 
This catastrophic collapse can be remedied by prepending special tokens (e.g. \texttt{\textbackslash n}), applying the chat template for instruct models (Table 8), or just not quantizing 0 v (\cite{yaqa}).
Regardless, \qts significantly improves upon \qss at all model sizes and bitrates, once again showing the dimensionality advantage of TCQ over VQ.

Table \ref{tab:llama31} shows results for Llama 3.1 instruct-tuned models, including Llama 3.1 405B.
At all sizes, QTIP achieves strong results.
Notably, QTIP is able to match or exceed PV-Tuning, a recent quantization method that focuses on better fine-tuning algorithms \cite{pvtuning}.
However, PV-Tuning is based off of AQLM and inherits its slow inference speed, making it significantly slower than QTIP.
Finally, Table \ref{tab:llama32} shows results for quantizing Llama 3.2 instruct-tuned models to 4 bits.
Since the embedding layers are very large relative to the decoder layers for small Llama 3 models ($\approx 500-750$MB), quantizing the decoder layers to fewer than 4 bits does not make a significant difference on the final model size. 
Here, QTIP is still able to achieve a meaningful end-to-end compression rate (2.5-3X) without degrading the final model. 

\subsection{Inference Speed}

Table \ref{tab:genspeed} shows the batch size 1 inference speed of \qt, \qs, and AQLM on Llama 2 7B and 70B with matrix fusion.
Here, the design choices of \qts and \qss become apparent.
Whereas AQLM uses a codebook that is too large to fit in cache and thus prevents fast inference, both \qts and \qss achieve significant speedups over FP16. 
Furthermore, while it is impressive that both \qss and \qts are $>2\times$ faster than AQLM, it is even more impressive that \qts is able to match \qs's throughput with an effective dimension size of 256, or $32 \times$ larger than \qs's.
This means that the improved quantization quality of \qts comes with \textit{no additional inference-time cost}.
Although our empirical throughput numbers were timed on NVIDIA GPUs, QTIP can be fast on a broad class of accelerators due to its flexibility.
QTIP only requires generating a pseudorandom Gaussian efficiently, and can work on devices with no cache as well as devices with lookup hardware. 
For example, if we were using a ARMv8 CPU, we could use the \texttt{vqtbl4q\_u8} NEON intrinsic to look up 16 indices in a 64-entry codebook. 
This would let us use a 6 bit 1D codebook with the HYB code (Q=6, V=1). 
Quantizing Llama 2 7B to 2 bits with this setup and w/out fine-tuning gives 6.89 Wikitext2 perplexity – essentially the same \sota quality as 3INST.